\chapter{数值级数}\label{ch:b1:series}

对无穷多个数求和，就是取其部分和的极限，
仅此而已。本章建立定义，并给出一年级可用的
收敛判别法：正项级数的比较法与等价法、
比值判别法、给出黎曼级数的积分比较法、
绝对收敛，以及交错级数定理。
更精细的理论（级数的乘积、分组求和、
函数项级数）
属于二年级。

\section{一般性质}

\begin{definition}\label{def:b1:series:def}
设 $(u_n)$ 为数列，称 $\sum u_n$ 为\emph{级数}\index{级数}，
它是由\emph{部分和} $S_N = \sum_{n=0}^{N} u_n$ 组成的数列。
若 $(S_N)$ 收敛，则称该级数\emph{收敛}，其极限即
为\emph{和} $\sum_{n=0}^{\infty} u_n$，并记 $R_N = \sum_{n > N} u_n =
S - S_N$ 为\emph{余项}，且趋于 $0$。
\end{definition}

\begin{example}[等比级数]\label{ex:b1:series:geometric}
对 $q \in \C$：$\;S_N = \sum_{n=0}^{N} q^n = \frac{1 -
q^{N+1}}{1-q}$（$q \neq 1$）。级数收敛当且仅当 $\abs q < 1$
（见\cref{exo:b1:seq:3}），且\index{等比级数}其和为
\[
\sum_{n=0}^{\infty} q^n = \frac{1}{1 - q} .
\]
\end{example}

\begin{example}[循环小数是等比级数]\label{ex:b1:series:periodicdecimal}
$0.363636\dots$ 是什么数？它的写法本身就是一个级数：
\[
0.\overline{36} = \sum_{k=1}^{\infty} \frac{36}{100^k}
= 36\cdot\frac{1/100}{1 - 1/100} = \frac{36}{99} =
\frac{4}{11} ,
\]
由等比求和，取 $q = \frac{1}{100}$。一般地，
循环节 $B$ 由 $p$ 位数字组成，无限重复时其值为
$\frac{B}{10^p - 1}$ --- 这正是周期性判别法
（见\cref{pb:b1:reals:1}）背后的机制，本章的语言
终于可以用一句话表述：小数展开是一个收敛
级数，当且仅当其和为有理数时，它最终是周期的。
第 10 章的数位机制仅凭上确界建立，
其实就是化了装的级数理论。
\end{example}

\begin{proposition}[基本性质]\label{prop:b1:series:first}
\begin{enumerate}
  \item 若 $\sum u_n$ 收敛，则 $u_n \to 0$。（逆命题
        \emph{不成立}：调和级数即是反例。）
  \item 线性性：收敛级数可逐项相加、数乘，
        且和随之相加、数乘。
  \item （裂项级数\index{裂项级数}）$\sum (v_{n+1} -
        v_n)$ 收敛当且仅当 $(v_n)$ 收敛，其和为 $\lim v_n -
        v_0$。
  \item 改变有限多项不影响收敛性
        （只改变和）。
\end{enumerate}
\end{proposition}

\begin{proof}
(1) $u_N = S_N - S_{N-1} \to S - S = 0$。调和级数有
$u_n = \frac1n \to 0$，但它发散（见\cref{exo:b1:seq:5}）。(2)
由极限的运算法则。(3) $S_N = v_{N+1} - v_0$。(4) 部分和
只相差一个最终恒定的量。
\end{proof}

\begin{example}[用等比余项规划项数]\label{ex:b1:series:georemainder}
当 $\abs q < 1$ 时，等比级数的余项
可以显式写出：
\[
R_N = \sum_{n = N+1}^{\infty} q^n = \frac{q^{N+1}}{1 - q} .
\]
这使精度要求在任何计算之前
就能换算为项数。例如计算 $\sum_{n\geq0} \bigl(\frac13\bigr)^n =
\frac32$，误差在 $10^{-10}$ 之内：需要 $\frac{(1/3)^{N+1}}{2/3} \leq
10^{-10}$，即 $3^{N} \geq \frac{3}{2}\cdot 10^{10}$，即 $N
\geq 22$（因 $3^{22} \approx 3.1\cdot10^{10}$），共需 23 项，
事先即知。凡周末问题中的等比速率估计，
（如以 $\frac13$ 为公比、用以求 $\ln 2$ 的级数，马钦的
反正切级数见\cref{pb:b1:taylor:1}），都是这套两行预算
穿上专业外衣后的样子。
\end{example}

\begin{example}[更长的裂项]\label{ex:b1:series:telescope3}
计算 $\sum_{n\geq1} \frac{1}{n(n+1)(n+2)}$。先作
部分分式分解（见\cref{ch:b1:fractions}）：
\[
\frac{1}{n(n+1)(n+2)}
= \frac{1/2}{n} - \frac{1}{n+1} + \frac{1/2}{n+2}
= \frac12\Bigl(\frac{1}{n(n+1)} - \frac{1}{(n+1)(n+2)}\Bigr),
\]
其中第二种形式 --- 即相邻两值之差 ---
$w_n = \frac{1}{n(n+1)}$ 的差 --- 正是裂项形式。于是
\[
\sum_{n=1}^{N} \frac{1}{n(n+1)(n+2)}
= \frac12\Bigl(w_1 - w_{N+1}\Bigr)
= \frac12\Bigl(\frac12 - \frac{1}{(N+1)(N+2)}\Bigr)
\longrightarrow \frac14 .
\]
收尾的要点：三项的部分分式直接写出时很少能
裂项相消；应先把它们改写为差 $w_n -
w_{n+1}$ 的形式 --- 收获的不仅是收敛，还有
精确的和，这是任何比较判别法都无法给出的。
\end{example}

\section{非负项级数}

\begin{theorem}[部分和有界]\label{thm:b1:series:positive}
若 $u_n \geq 0$ 对一切 $n$ 成立，则部分和单调递增，故：$\sum
u_n$ 收敛 $\iff$ 其部分和有上界。由此得到
\emph{比较判别法}：若 $0 \leq u_n \leq v_n$ 对
所有（充分大的）$n$ 成立，
\[
\sum v_n \text{ 收敛} \implies \sum u_n \text{ 收敛},
\qquad
\sum u_n \text{ 发散} \implies \sum v_n \text{ 发散}.
\]
又有\emph{等价判别法}：若 $u_n \sim v_n$ 且 $v_n \geq
0$，两级数同敛散。
\end{theorem}

\begin{proof}
由单调极限定理（\cref{thm:b1:seq:monotone}）得第一
点，由部分和的比较得第二点。等价关系：对
充分大的 $n$，$\frac12 v_n \leq u_n \leq 2 v_n$（由 $\sim$ 的定义，
取 $\varepsilon = \frac12$），两个方向的比较判别法均适用。
\end{proof}

\begin{example}[用等价量证明发散]\label{ex:b1:series:equivdiverge}
$\sum_{n\geq1} n\sin\dfrac{1}{n^2}$ 的敛散性如何？由于
$\frac{1}{n^2} \to 0$，且 $\sin h \sim h$（在 $0$ 附近）：
\[
n\sin\frac{1}{n^2} \;\sim\; n\cdot\frac{1}{n^2} = \frac1n ,
\]
由等价判别法，调和级数的发散性被传递：
级数发散 --- 尽管通项趋于
$0$。一次展开，一个尺度，一个结论；同样的两步
套路（先作等价，再查黎曼级数或几何级数）可判定
\cref{exo:b1:series:3}中的全部四个级数。
\end{example}

\begin{example}[一行内的等价判别法]\label{ex:b1:series:equivtest}
$\sum_{n \geq 1} \frac{\sqrt{n+1} - \sqrt n}{n}$ 的敛散性如何？
分子有理化：
\[
\frac{\sqrt{n+1} - \sqrt n}{n}
= \frac{1}{n\,(\sqrt{n+1} + \sqrt n)}
\sim \frac{1}{2\,n^{3/2}} ,
\]
这是收敛的黎曼尺度（$\alpha = \frac32 > 1$）：级数
收敛。整个判断只用了一次等价和一次查表
--- 前提是通项非负，本题确实如此。结论性的
看法是：对于正项级数，整个收敛
理论就是一本\emph{尺度词典}（$n^{-\alpha}$、$q^n$、
$\frac{1}{n(\ln n)^\alpha}$），以及用等价量替换通项的
许可。分析的工作在于渐近行为（
见\cref{ch:b1:taylor}），绝不在求和之中。
\end{example}

\begin{theorem}[积分比较；黎曼级数]\label{thm:b1:series:riemann}
设 $f$ 为连续、非负且\emph{递减}的函数，定义在
$\intco{1}{+\infty}$ 上。则
\[
\int_1^{N+1} f(t)\,\dd t \;\leq\; \sum_{n=1}^{N} f(n) \;\leq\; f(1)
+ \int_1^{N} f(t)\,\dd t ,
\]
故 $\sum f(n)$ 收敛当且仅当 $\bigl(\int_1^x f\bigr)$ 有界。特别地，
对 $\alpha \in \R$：\index{黎曼级数}
\[
\sum_{n \geq 1} \frac{1}{n^\alpha} \text{ 收敛}
\iff \alpha > 1,
\]
且 $\sum_{n=1}^{N} \frac1n = \ln N + O(1)$。
\end{theorem}

\begin{proof}
当 $n \leq t \leq n+1$ 时，由单调性得 $f(n+1) \leq f(t)
\leq f(n)$。在 $\intcc{n}{n+1}$（长度为
$1$ 的区间）上积分，得：
\[
f(n+1) \;\leq\; \int_n^{n+1} f(t)\,\dd t \;\leq\; f(n) .
\]
对 $n = 1, \dots, N-1$ 求和右侧不等式，得
$\int_1^{N} f \leq \sum_{n=1}^{N-1} f(n)$，再加上
$f(N) \leq f(1)$，即得上界估计；对
$n = 1, \dots, N$ 求和左侧不等式，得 $\sum_{n=2}^{N+1} f(n) \leq
\int_1^{N+1} f$，经指标替换后即为下界估计。
收敛性：部分和与积分 $\int_1^x f$
相差不超过常数 $f(1)$，且二者都
单调不减，故其一有界当且仅当其二有界
（见\cref{thm:b1:series:positive}）。对 $f(t)
= t^{-\alpha}$ ($\alpha \neq 1$): $\int_1^x t^{-\alpha}\dd t =
\frac{x^{1-\alpha} - 1}{1 - \alpha}$，有界当且仅当 $\alpha > 1$；对
$\alpha = 1$ 时积分为 $\ln x \to \infty$，而夹逼
估计给出 $\ln(N+1) \leq H_N \leq 1 + \ln N$。对 $\alpha \leq 0$，
通项不趋于 $0$。
\end{proof}

\begin{example}[调和级数的累加]\label{ex:b1:series:harmonicstack}
调和级数要累加多少项，才能超过 $20$？
夹逼估计 $\ln(N+1) \leq H_N \leq 1 + \ln N$ 无需任何求和即可
回答：$H_N \geq 20$ 要求 $1 + \ln N \geq 20$，
即 $N \geq \eu^{19} \approx 1.8\cdot10^{8}$，而
一旦 $\ln(N + 1) \geq 20$ 即可保证，此时 $N \approx \eu^{20} \approx
4.9\cdot10^{8}$。（周末问题将此估计改进为 $N \approx
\eu^{20 - \gamma} \approx 2.7\cdot10^{8}$，借助欧拉常数。）
结论性的看法是：积分比较不仅
判定敛散性 --- 它能\emph{定位}部分和，
以对数精度把原本无望的计算
（数以亿计的项）化为两行估计。
\end{example}

\begin{theorem}[比值判别法（达朗贝尔）]\label{thm:b1:series:ratio}
设 $u_n > 0$ 且 $\frac{u_{n+1}}{u_n} \to \ell$。\index{比值判别法}
\begin{itemize}
  \item 若 $\ell < 1$：$\sum u_n$ 收敛；
  \item 若 $\ell > 1$：$u_n \to +\infty$，发散；
  \item 若 $\ell = 1$：无结论（$\sum \frac1n$ 发散，$\sum
        \frac{1}{n^2}$ 收敛）。
\end{itemize}
\end{theorem}

\begin{proof}
若 $\ell < 1$，取 $q \in \intoo{\ell}{1}$，则自某项 $N$ 之后，
$u_{n+1} \leq q\,u_n$，由归纳法得 $u_n \leq u_N q^{\,n-N}$：
与等比级数作比较。若 $\ell > 1$，则自某项 $N$ 之后
数列 $(u_n)$ 递增，故不可能趋于 $0$
（其极限若存在，则 $\geq u_N > 0$）；由
\cref{prop:b1:series:first}~(1)得发散 --- 事实上 $u_n
\geq u_N q^{n-N}$，其中 $q > 1$，故 $u_n \to \infty$。
\end{proof}

\begin{example}\label{ex:b1:series:ratio}
$\sum \frac{x^n}{n!}$ 对每个 $x > 0$ 收敛：比值
$\frac{x}{n+1} \to 0$。其和为 $\eu^x$：由泰勒–拉格朗日
公式（\cref{thm:b1:taylor:lagrange}）在 $\intcc{0}{x}$ 上，有
\[
\Bigl| \eu^x - \sum_{k=0}^{n} \frac{x^k}{k!} \Bigr|
\leq \eu^{x}\, \frac{x^{n+1}}{(n+1)!} \xrightarrow[n\to\infty]{} 0 ,
\]
该界趋于 $0$，因阶乘占主导
（\cref{exo:b1:integration:9}~(1) 用到了同一事实）。同样的
论证对 $\sin$、$\cos$、$\sinh$、$\cosh$ 的级数，
在整个 $\R$ 上同样可以求和。
\end{example}

\begin{example}[比值判别法是充分的，而非必要的]\label{ex:b1:series:rationotnecessary}
设 $u_n = 2^{-n}$（$n$ 为偶数），$u_n = 2^{-n-2}$（
$n$ 为奇数）。相邻比值在 $\frac{1}{8}$ 与
$\frac12\cdot4 = 2$，故 $\frac{u_{n+1}}{u_n}$ 无极限，
达朗贝尔判别法失效——然而 $u_n \leq 2^{-n}$，而比较
判别法即刻判定其收敛。该判别法的前提（比值
\emph{收敛}）是一个真实的限制：它适用于
具有单一主导乘法结构（阶乘、幂）的项，
而对起伏不定的项则失效。比值表现失常时，
退回与几何型包络的比较——
这恰是比值判别法的本质，其证明即可看出。
\end{example}

\begin{example}[阶乘对决中的比值判别法]\label{ex:b1:series:ratiobinom}
判定 $\sum_{n\geq0} \dfrac{(n!)^2}{(2n)!}$ 的敛散性（它是
中心二项式系数的倒数，相差因子 $n + 1$）？
比值可消去阶乘：
\[
\frac{u_{n+1}}{u_n}
= \frac{((n+1)!)^2}{(n!)^2}\cdot\frac{(2n)!}{(2n+2)!}
= \frac{(n+1)^2}{(2n+1)(2n+2)}
\longrightarrow \frac14 < 1 :
\]
收敛，且尚有余裕——各项本质上按
$4^{-n}$ 的速度衰减，与 $\binom{2n}{n} \geq
\frac{4^n}{2n+1}$ 相符（见\cref{pb:b1:integration:1}）。
要点是：阶乘之商恰是比值判别法
所能处理的对象——每个阶乘皆消去，化为关于
$n$ 的有理函数，其极限可由首项读出。
\end{example}

\section{绝对收敛；交错级数}

\begin{theorem}[绝对收敛]\label{thm:b1:series:absolute}
若 $\sum \abs{u_n}$ 收敛（\emph{绝对
收敛}\index{绝对收敛}），则 $\sum u_n$ 收敛，
且 $\bigl|\sum u_n\bigr| \leq \sum \abs{u_n}$。此结论对
实项与复项均成立。
\end{theorem}

\begin{proof}
对 $M > N$，部分和满足（柯西准则，
见\cref{thm:b1:seq:complete}）：
\[
\abs{S_M - S_N} = \Bigl| \sum_{n=N+1}^{M} u_n \Bigr|
\leq \sum_{n=N+1}^{M} \abs{u_n},
\]
当 $N$ 充分大时该式很小，因为 $\sum\abs{u_n}$ 的部分和
构成柯西序列。故 $(S_N)$ 亦为柯西序列，从而收敛。
该不等式由有限和的三角不等式取极限得到。
\end{proof}

\begin{example}[绝对收敛：实数与复数情形]\label{ex:b1:series:absexamples}
$\sum_{n\geq1} \frac{\sin n}{n^2}$：各项符号
毫无规律地变化（事实上 $(\sin n)$ 在 $\intcc{-1}{1}$ 中稠密，
见\cref{exo:b1:seq:12}），且看不出任何
交错结构。绝对收敛一举解决了全部问题：
$\bigl|\frac{\sin n}{n^2}\bigr| \leq \frac{1}{n^2}$，
这是收敛的标尺，故级数收敛。同一护盾
在 $\C$ 上同样有效：$\sum_{n\geq1}\frac{\eu^{\iu n}}{n^2}$ 收敛，
因为 $\bigl|\frac{\eu^{\iu n}}{n^2}\bigr| = \frac{1}{n^2}$
——符号的分布，哪怕是二维的，也都无关紧要
一旦模的级数收敛。要点在于：绝对收敛
是本章唯一从不追问符号
如何排布的工具，应首先尝试它，
（见\cref{met:b1:series:decide}），而将精细的判别法
留给那些不满足绝对收敛的级数。
\end{example}

\begin{theorem}[交错级数判别法]\label{thm:b1:series:alternating}
设 $(a_n)$ 单调递减，且 $a_n \to 0$。则交错
级数 $\sum (-1)^n a_n$ 收敛；其和位于任意两个
相邻部分和之间，且\index{交错级数}
\[
\abs{R_N} = \Bigl| \sum_{n > N} (-1)^n a_n \Bigr| \leq a_{N+1} .
\]
\end{theorem}

\begin{proof}
偶数与奇数部分和相邻：$S_{2p+2} - S_{2p} =
a_{2p+2} - a_{2p+1} \leq 0$（递减），$S_{2p+1} - S_{2p-1} =
a_{2p} - a_{2p+1} \geq 0$（递增），且 $S_{2p} - S_{2p+1} =
a_{2p+1} \to 0$。由\cref{thm:b1:seq:adjacent}，二者有公共极限
$S$，由两子列准则
（\cref{prop:b1:seq:subsequences}）知它即为 $(S_N)$ 的极限；
此外 $S$ 夹在相邻部分和之间，且
$\abs{S - S_N}$ 不超过与下一部分和的间距 $a_{N+1}$。
\end{proof}

\begin{example}[交错调和级数]\label{ex:b1:series:ln2}
$\sum_{n \geq 1} \frac{(-1)^{n-1}}{n}$ 收敛（由交错级数判别法），
但非绝对收敛（调和级数发散）。其和为 $\ln 2$，由
有限几何恒等式 $\frac{1}{1+t} = \sum_{k=0}^{n-1} (-t)^k +
\frac{(-t)^n}{1+t}$，在 $\intcc{0}{1}$ 上积分，得：
\[
\ln 2 = \sum_{k=1}^{n} \frac{(-1)^{k-1}}{k}
+ (-1)^n \int_0^1 \frac{t^n}{1+t}\,\dd t ,
\qquad
0 \leq \int_0^1 \frac{t^n}{1+t}\,\dd t \leq \frac{1}{n+1} \to 0 .
\]
收敛极其缓慢（$R_N \approx \frac{1}{N}$）——
交错级数依靠正负相消而收敛，而非依靠各项之小。
\end{example}

\begin{omfigure}
\begin{tikzpicture}[x=0.78cm, y=7.2cm]
  \draw[-Stealth] (0.3,0.42) -- (11,0.42)
    node[below, font=\small] {$N$};
  \draw (1,0.415) -- (1,0.425) node[below=1pt, font=\small]
    {$1$};
  \draw (5,0.415) -- (5,0.425) node[below=1pt, font=\small]
    {$5$};
  \draw (10,0.415) -- (10,0.425) node[below=1pt, font=\small]
    {$10$};
  \draw[gray, dashed] (0.3,0.6931) -- (10.7,0.6931)
    node[right, font=\small, text=black] {$\ln 2$};
  \node[left, font=\small] at (0.3,1) {$1$};
  \node[left, font=\small] at (0.3,0.5) {$0.5$};
  \draw[omDef, thick]
    plot coordinates {(1,1) (2,0.5) (3,0.8333) (4,0.5833)
    (5,0.7833) (6,0.6167) (7,0.7595) (8,0.6345)
    (9,0.7456) (10,0.6456)};
  \fill[omDef] (1,1) circle (1.5pt);
  \fill[omDef] (2,0.5) circle (1.5pt);
  \fill[omDef] (3,0.8333) circle (1.5pt);
  \fill[omDef] (4,0.5833) circle (1.5pt);
  \fill[omDef] (5,0.7833) circle (1.5pt);
  \fill[omDef] (6,0.6167) circle (1.5pt);
  \fill[omDef] (7,0.7595) circle (1.5pt);
  \fill[omDef] (8,0.6345) circle (1.5pt);
  \fill[omDef] (9,0.7456) circle (1.5pt);
  \fill[omDef] (10,0.6456) circle (1.5pt);
\end{tikzpicture}

{\small 部分和 $S_N$ 来自交错调和级数 $1 - \frac12 + \frac13 - \cdots$，
在每一步都越过其极限
$\ln 2$：奇数项部分和在上方，偶数项部分和在下方，
每次跳跃的幅度为 $\frac{1}{N+1}$。夹逼
正是\cref{thm:b1:series:alternating}证明的直观呈现 --- 而
收窄过程极其缓慢（$\abs{S_N - \ln 2} \approx
\frac{1}{2N}$，见\cref{pb:b1:series:1}），这也解释了为何
无人用这种方法计算 $\ln 2$。}
\end{omfigure}

\begin{remark}[级数的常见误区]
（i）\emph{等价判别法须顾及符号}：设 $v_n =
\frac{(-1)^n}{\sqrt n}$ 及 $u_n = v_n + \frac1n$，则
$\frac{u_n}{v_n} = 1 + \frac{(-1)^n}{\sqrt n} \to 1$，故 $u_n
\sim v_n$；然而 $\sum v_n$ 依交错级数判别法收敛，而
$\sum u_n = \sum v_n + \sum \frac1n$ 发散。等价性控制的是
项的\emph{大小}，而对带符号的级数，大小并
不决定一切 --- 该判别法的叙述与正确性，仅对
（自某项起）非负的项成立。（ii）\emph{$u_n \to 0$
不能说明任何问题}：调和级数是永恒的反例；
其逆向陈述
（\cref{prop:b1:series:first}~(1)）仅是一个快速的
发散判别法。（iii）\emph{比值极限为 $1$ 无法说明
收敛}：$\sum\frac1n$ 与 $\sum\frac{1}{n^2}$ 的比值都
$\to 1$；应改用黎曼尺度或积分比较。
（iv）\emph{交错级数须单调递减}：$\sum
\frac{(-1)^n}{n + (-1)^n}$ 看似交错，却只能
借助展开来处理（见\cref{exo:b1:series:5}）；
\cref{ch:b1:taylor}中的周末问题（第 23 题）表明，该判别法
在缺少单调性时可能彻底失效。（v）\emph{分组与重排
并非无代价}：对收敛级数加括号无害，
但可能从发散中制造出收敛
（如 $1 - 1 + 1 - \cdots$ 两两分组），而重排
则可以改变级数本身的和 --- 这一点在本章的
周末问题（\cref{pb:b1:series:1}）中有所上演。
\end{remark}

\begin{method}[判定级数的敛散性]\label{met:b1:series:decide}
\begin{enumerate}
  \item $u_n \to 0$ 是否成立？若否，则级数发散，停止。
  \item 若各项非负：求 $u_n$ 的等价量（展开式，
        见\cref{ch:b1:taylor}！），与黎曼尺度或几何尺度比较；
        阶乘与幂次宜用比值判别法；
        单调递减的 $f(n)$ 宜用积分比较。
  \item 符号可变：先考察绝对收敛；若不成立，则
        交错级数判别法（须仔细检验\emph{单调递减}）；此外
        则需借助本科二年级的工具。
\end{enumerate}
\end{method}

\begin{example}[奇数分母：半截裂项相消]\label{ex:b1:series:oddtelescope}
求 $\sum_{n \geq 1} \dfrac{1}{4n^2 - 1}$。
部分分式：$\frac{1}{(2n-1)(2n+1)} =
\frac12\bigl(\frac{1}{2n-1} - \frac{1}{2n+1}\bigr)$，于是
\[
\sum_{n=1}^{N} \frac{1}{4n^2 - 1}
= \frac12\Bigl(1 - \frac{1}{2N+1}\Bigr)
\longrightarrow \frac12 .
\]
比较 $\sum \frac{1}{n(n+1)} = 1$
（见\cref{exo:b1:series:1}）：二者有相同的裂项骨架，但
此处的相邻项在奇数序列中相隔两步，而
因子 $\frac12$ 记录了这一步长。末尾点明：
裂项相消是一种视角的转换，而非技巧 --- 只要
通项为某有极限的数列的差 $w_n - w_{n+1}$，
则和为 $w_1 - \lim w$，恰为
\cref{prop:b1:series:first}~(3)。
\end{example}

\begin{remark}[回望分析学的主线]
本章是本分册分析内容汇聚之处，每个判别法
都指明了它的源头。部分和有界，即是单调极限定理，
（见\cref{ch:b1:seq}），其本身是
在\cref{ch:b1:reals}中的完备性公理；绝对收敛是柯西
准则；积分判别法来自\cref{ch:b1:integration}中对面积的
刻画；通项的等价量即是
\cref{ch:b1:taylor}中的展开；而交错级数判别法是
相邻数列引理换了周末装束后的形式。倒过来读，
这条主线解释了每一章究竟是
\emph{为了什么} --- 而贯穿其中的周末问题
（$b$ 进数字、切萨罗–施托尔茨、无理性机器、
欧拉常数）都是同样的几个思想，在越来越高的层次上相会。
此后的线性代数改变的是对象，
而非标准：以经证明的误差界来精确陈述的习惯，
从极限过渡到维数的时候依然保留。
\end{remark}

\begin{remark}[级数的下一站]
本章为本册的分析收尾，并打开三扇
门。在本科二年级分册中，级数获得了一个变量（$\sum
a_n x^n$：幂级数及其收敛半径），然后
发展出以函数为对象的理论（傅里叶级数）；绝对收敛与条件收敛的二分，
在下面的周末问题中得到生动体现，
将成为二者共同的基石。在概率论
（本科三年级分册）中，离散随机变量的期望
\emph{就是}级数，而绝对收敛正是
使它们得以良定义的条件。而黎曼级数 $\sum n^{-s}$，
推广到复数 $s$ 时，便成为 ζ 函数 --- 数学中被研究得最多的
一个级数。
\end{remark}

\section{练习}

\begin{exercise}[$\star$]\label{exo:b1:series:1}
判断下列级数的敛散性（若可裂项相消，则求其和）：
\[
\sum_{n\geq1} \frac{1}{n(n+1)},
\qquad
\sum_{n\geq2} \ln\Bigl(1 - \frac{1}{n^2}\Bigr),
\qquad
\sum_{n\geq0} \frac{3^n + 4^n}{5^n} .
\]
\end{exercise}

\begin{exercise}[$\star$]\label{exo:b1:series:2}
判断下列级数的敛散性：$\;\sum \dfrac{n^2}{2^n}$；$\;\sum \dfrac{n!}{n^n}$；
$\;\sum \dfrac{2^n\,n!}{n^n}$；$\;\sum \dfrac{3^n\,n!}{n^n}$。
\emph{（比值判别法；回忆 $\bigl(1 + \frac1n\bigr)^n \to \eu$。）}
\end{exercise}

\begin{exercise}[$\star$]\label{exo:b1:series:3}
判断下列级数的敛散性：$\;\sum \sin\dfrac{1}{n^2}$；$\;\sum
\Bigl(1 - \cos\dfrac1n\Bigr)$；$\;\sum \dfrac{1}{\sqrt{n(n+1)}}$；
$\;\sum \dfrac{\ln n}{n^2}$ \emph{（与 $n^{-3/2}$ 比较）}。
\end{exercise}

\begin{exercise}[$\star$]\label{exo:b1:series:4}
证明 $\sum_{n\geq1} \frac{1}{n^2}$ 收敛，且其和 $\leq 2$，
利用 $\frac{1}{n^2} \leq \frac{1}{n(n-1)}$（$n \geq 2$），并
借助裂项相消给出上界。
\end{exercise}

\begin{exercise}[$\star\star$]\label{exo:b1:series:5}
判断下列级数的敛散性：$\;\sum \dfrac{(-1)^n}{\sqrt n}$，$\;\sum
\dfrac{(-1)^n}{n + (-1)^n}$ \emph{（先展开：交错级数判别法
不能直接应用 --- 为什么？）}，以及 $\;\sum
\sin\bigl(\pi\sqrt{n^2+1}\,\bigr)$ \emph{（模 $\pi$ 约化：
$\sqrt{n^2+1} = n + \frac{1}{2n} + O(n^{-3})$）}。
\end{exercise}

\begin{exercise}[$\star\star$]\label{exo:b1:series:6}
对哪些 $\alpha > 0$，级数 $\sum_{n \geq 2}
\dfrac{1}{n (\ln n)^{\alpha}}$ 收敛？\emph{（积分比较法；
作代换 $u = \ln t$。）}
\end{exercise}

\begin{exercise}[$\star\star$]\label{exo:b1:series:7}
设 $u_n = \dfrac{1}{n} - \ln\Bigl(1 + \dfrac1n\Bigr)$。证明
$0 \leq u_n \leq \dfrac{1}{2n^2}$，$\sum u_n$ 收敛，并
由此推出欧拉常数\index{欧拉常数}的存在性：
\[
\gamma = \lim_{N \to \infty} \Bigl( \sum_{n=1}^{N} \frac 1n - \ln N
\Bigr) .
\]
\end{exercise}

\begin{exercise}[$\star\star$]\label{exo:b1:series:8}
求下列级数的和：
\[
\sum_{n=1}^{\infty} \frac{1}{n(n+2)}
\qquad\text{和}\qquad
\sum_{n=0}^{\infty} \frac{n}{2^n} .
\]
\emph{（第一题用部分分式。第二题：先求
$\sum_{n=1}^{N} n x^{n-1}$ 的闭式，再令 $N \to \infty$，
在 $x = \frac12$ 处。）}
\end{exercise}

\begin{exercise}[$\star\star\star$]\label{exo:b1:series:9}
（柯西凝聚判别法）设 $(u_n)$ 非负且单调递减。
证明
\[
\sum_{n \geq 1} u_n \text{ 收敛}
\iff
\sum_{k \geq 0} 2^k\, u_{2^k} \text{ 收敛},
\]
通过比较相邻 $2$ 的幂之间的各组项。
由此推出黎曼判别法，以及
\cref{exo:b1:series:6}。
\end{exercise}

\begin{exercise}[$\star\star\star$]\label{exo:b1:series:10}
利用\cref{ex:b1:series:ln2}中的积分恒等式，
将其改用于 $\frac{1}{1+t^2}$，证明莱布尼茨公式
\[
\frac{\pi}{4} = \sum_{n=0}^{\infty} \frac{(-1)^n}{2n+1}
= 1 - \frac13 + \frac15 - \frac17 + \cdots
\]
并给出误差界 $\abs{R_N} \leq \frac{1}{2N+3}$。
\end{exercise}

\begin{exercise}[$\star\star$]\label{exo:b1:series:11}
判断 $\displaystyle\sum_{n \geq 1} \frac{1}{n^{1 + 1/n}}$ 的敛散性。
\emph{（求 $n^{1/n}$ 的极限，并求通项的等价量：
黎曼判别法需要一个\emph{固定}的
指数。）}
\end{exercise}

\begin{exercise}[$\star\star\star$]\label{exo:b1:series:12}
设 $(u_n)$ 为非负且\emph{递减}的数列，级数 $\sum u_n$
收敛。证明 $n\,u_n \to 0$ \emph{（将 $n\,u_{2n}$
用段和 $\sum_{k=n+1}^{2n} u_k$ 控制，并利用柯西
准则）}。证明逆命题不成立，并说明
单调性假设不可去掉。
\end{exercise}

\section{周末问题：欧拉常数与改变其和的
级数}

\begin{problem}[{周末问题 --- $H_n = \ln n + \gamma +
\frac{1}{2n} + O(n^{-2})$，并将 $1 - \frac12 +
\frac13 - \dots$ 重排为 $\frac{\ln 2}{2}$}]
\label{pb:b1:series:1}
两个故事共享调和级数。其一是其发散性的
精确核算：$H_n - \ln n$ 收敛于欧拉
常数 $\gamma$（见\cref{exo:b1:series:7}），本题
将该结论加细为双侧估计 $\frac{1}{2(n+1)}
\leq H_n - \ln n - \gamma \leq \frac{1}{2n}$，并以手算验证 $\gamma
= 0.5772\dots$。其二是条件收敛的反常现象：
交错调和级数之和为 $\ln 2$
（见\cref{ex:b1:series:ln2}），而\emph{同样的项，次序不同}，
其和为 $\frac{\ln 2}{2}$ ——或 $\ln 2 +
\frac12\ln\frac pq$（对任意 $p, q$），或为任一实数
（黎曼）。两个故事本是一回事：重排后的和
均\emph{借助} $\gamma$ 律计算。\index{欧拉常数}
\index{级数的重排}

\medskip
\noindent\textbf{第一部分 --- $\gamma$ 的夹逼。}设 $a_n = H_n
- \ln n$ 和 $b_n = H_n - \ln(n+1)$。
\begin{enumerate}
  \item 利用 $\frac{t}{1+t} \leq \ln(1 + t) \leq t$，证明
        $(a_n)$ 递减，$(b_n)$ 递增，且二者
        相邻；其公共极限为 $\gamma$，并有 $b_n \leq
        \gamma \leq a_n$ 对每个 $n$ 成立。
  \item 数值初试：由 $H_{10} = 2.928968\dots$，
        将 $\gamma$ 夹在 $b_{10} = 0.5311$ 与 $a_{10} =
        0.6264$。此粗糙夹逼若要精确到四位小数，
        需要多大的 $n$？
  \item 证明精确的余项表示 $a_n - \gamma =
        \sum_{k \geq n} w_k$（部分和的极限），其中
        \[
        w_k = a_k - a_{k+1}
        = \ln\Bigl(1 + \frac1k\Bigr) - \frac{1}{k+1}
        = \int_k^{k+1} \frac{(k + 1 - t)}{t\,(k+1)}\,\dd t ,
        \]
        并由积分形式推出双侧估计
        $\dfrac{1}{2(k+1)^2} \leq w_k \leq \dfrac{1}{2k(k+1)}$。
\end{enumerate}

\medskip
\noindent\textbf{第二部分 --- $\frac{1}{2n}$ 律。}
\begin{enumerate}[resume]
  \item 将第 3 问的界相加（两侧皆可裂项相消，或
        与裂项和比较），从而得到该律：
        \[
        \frac{1}{2(n+1)} \;\leq\; H_n - \ln n - \gamma \;\leq\;
        \frac{1}{2n} \qquad (n \geq 1).
        \]
  \item 推出 $H_n = \ln n + \gamma + \frac{1}{2n} +
        O\bigl(\frac{1}{n^2}\bigr)$；更精确地，证明
        $\gamma_n = H_n - \ln n - \frac{1}{2n}$ 满足
        $-\frac{1}{2n(n+1)} \leq \gamma_n - \gamma \leq 0$。
  \item 取 $n = 100$，验证四位小数：已知 $H_{100} =
        5.1873775\dots$，计算 $\gamma_{100} = 0.577207\dots$
        并得出 $\gamma = 0.5772 \pm 5\cdot10^{-5}$（真值
        为 $0.5772156\dots$）。
  \item 该律的两个推论，后文均会用到：当 $m \to
        \infty$ 时，
        \[
        H_{2m} - H_m = \ln 2 - \frac{1}{4m} +
        O\Bigl(\frac{1}{m^2}\Bigr),
        \qquad
        \sum_{j=1}^{m} \frac{1}{2j-1} = \frac{\ln m}{2} + \ln 2
        + \frac\gamma2 + o(1) ,
        \]
        后者借助 $\sum_{j \leq m} \frac{1}{2j-1} = H_{2m}
        - \frac12 H_m$，同理 $\sum_{j=1}^{m}
        \frac{1}{2j} = \frac{\ln m}{2} + \frac\gamma2 + o(1)$。
\end{enumerate}

\medskip
\noindent\textbf{第三部分 --- 交错调和级数，
精确到二阶。}
\begin{enumerate}[resume]
  \item 用归纳法或分组法证明恒等式
        $\sum_{k=1}^{2m} \frac{(-1)^{k-1}}{k} = H_{2m} - H_m$，
        由此推出其和 $\ln 2$（再次得到），以及精确的
        收敛速度：
        \[
        \sum_{k=1}^{2m} \frac{(-1)^{k-1}}{k}
        = \ln 2 - \frac{1}{4m} + O\Bigl(\frac{1}{m^2}\Bigr) .
        \]
  \item 推出交错调和级数的渐近误差
        在\emph{任意}指标处 $S - S_N \sim
        \frac{(-1)^N}{2N}$ ——比
        最坏情形界 $a_{N+1} \approx \frac1N$ 小一半（见
        \cref{thm:b1:series:alternating}）。
  \item （免费的加速）证明平均部分和
        $\tilde S_N = \frac{S_N + S_{N+1}}{2}$ 满足 $\tilde
        S_N = \ln 2 + O\bigl(\frac{1}{N^2}\bigr)$。验证：$S_{10}
        = 0.64563$, $S_{11} = 0.73654$, $\tilde S_{10} =
        0.69109$，对照 $\ln 2 = 0.69315$：一次平均
        换得两位小数的精度。
  \item 用两句话说明，为何此类技巧无法帮助\emph{正}的
        发散尾项现象（如第 2 问的夹逼）：交错误差
        振荡（符号
        $(-1)^N$），故平均会消去其主项，而
        $\gamma$ 夹逼的误差 $\frac{1}{2n}$ 符号
        恒定。（对 $a_n$ 与 $b_n$ 取平均\emph{确有}帮助：
        将 $\frac{a_n + b_n}{2}$ 与中点估计
        $H_n - \ln\bigl(n + \frac12\bigr)$ 相联系，并证其误差为
        $O\bigl(\frac{1}{n^2}\bigr)$。）
\end{enumerate}

\medskip
\noindent\textbf{第 IV 部分 --- 刚性及其失效。}
\begin{enumerate}[resume]
  \item 证明交错调和级数的正项部分 $\sum \frac{1}{2j-1}$ 与
        负项部分 $\sum \frac{1}{2j}$ 均发散 ---
        这正是
        \emph{条件收敛}的标志。
  \item 证明第 12 问背后的一般结论：若
        $\sum u_n$ 收敛而 $\sum \abs{u_n}$ 发散，则
        正项部分构成的级数 $\sum u_n^+$ 与
        负项部分构成的级数 $\sum u_n^-$ 均发散 \emph{（由
        $u_n^\pm = \frac{\abs{u_n} \pm u_n}{2}$：若其中一个收敛，
        则另一个亦收敛，从而 $\sum\abs{u_n}$ 收敛）}。
        这一取之不尽的正负项“质量”储备，
        正是黎曼配方将要耗用的。
  \item （刚性）证明：若 $\sum u_n$ \emph{绝对}收敛，
        且 $\sigma \colon \N \to \N$ 为
        双射，则 $\sum u_{\sigma(n)}$ 收敛于同一个和
        \emph{（取 $N$ 充分大时，重排后的前 $M$ 项
        包含 $u_0, \dots, u_N$；比较部分和
        借助尾项 $\sum_{n > N}\abs{u_n}$）}。
  \item （黎曼配方）设 $t \in \R$。描述交错调和级数的贪心重排：
        即先依次取正项 $1, \frac13, \frac15, \dots$，
        直到部分和首次超过 $t$，再依次取负项 $-\frac12, -\frac14, \dots$，
        直到部分和首次低于 $t$，并如此反复。
        证明每一项恰好被使用一次，
        且首次穿越之后，部分和
        与 $t$ 的距离始终不超过最后所用项的绝对值，
        由此得出重排级数收敛于 $t$：
        \emph{任意}预先给定的和
        均可达到。
\end{enumerate}

\medskip
\noindent\textbf{第 V 部分 --- $(p, q)$ 公式。}取定整数
$p, q \geq 1$。将交错调和级数分块重排：
每块先取 $p$ 个正项（即接下来的奇数倒数），再取 $q$ 个负项
（即接下来的偶数倒数），并如此重复。
\begin{enumerate}[resume]
  \item 验证这确是一个真正的重排（每一项恰好出现一次），
        并说明当 $(p, q) = (1, 2)$ 时，它的形式为
        \[
        1 - \frac12 - \frac14 + \frac13 - \frac16 - \frac18 +
        \frac15 - \frac1{10} - \frac1{12} + \dots
        \]
  \item （精确减半）对 $(p, q) = (1, 2)$，证明
        分块恒等式
        \[
        \frac{1}{2k-1} - \frac{1}{4k-2} - \frac{1}{4k}
        = \frac12\Bigl(\frac{1}{2k-1} - \frac{1}{2k}\Bigr),
        \]
        并由此推出 \emph{精确}关系 $T_{3K} = \frac12
        S_{2K}$，即重排与原来的部分和之间的关系：
        和的减半在每个有限阶段都清晰可见，
        而不只是在极限中。
  \item 证明经过 $K$ 个完整块之后的部分和
        等于 $\sum_{j=1}^{pK} \frac{1}{2j-1} -
        \sum_{j=1}^{qK} \frac{1}{2j}$，并借助第 7 问求其极限，
        结果为
        \[
        \ln 2 + \frac12 \ln\frac pq .
        \]
  \item 控制\emph{块内}的部分和（各项趋于 $0$），
        并得出按 $(p,
        q)$ 重排的级数收敛于 $\ln 2 + \frac12\ln
        \frac pq$。特别地，$(1, 2)$ 给出 $\frac{\ln
        2}{2}$：用前九项加以验证，$T_9 =
        0.3083$，并逐步趋近于 $0.3466$。
  \item 检验与范围：$(1,1)$ 给出 $\ln 2$；
        $(2,1)$ 给出 $\frac32\ln 2$；哪些和可由 $(p, q)$ 块达到，
        这一可数的“菜单”
        与黎曼的完整“菜单”（第 14 问）相比如何？
\end{enumerate}

\medskip
\noindent\textbf{第 VI 部分 --- 尾声：$\gamma$ 的应用与
综合。}
\begin{enumerate}[resume]
  \item 确定收敛级数
        $\sum_{k\geq1} \bigl(\frac1k - \ln\frac{k+1}{k}\bigr)$
        的和（见\cref{exo:b1:series:7}），并证明它等于 $\gamma$。
  \item 对目标值 $t =
        1$，运行黎曼配方（第 14 问），列出所产生的前十二项（$1, \frac13,
        -\frac12, \frac15, -\frac14, \frac17, \frac19, -\frac16,
        \frac1{11}, \frac1{13}, -\frac18, \frac1{15}$),
        并计算部分和（$\approx 0.980$） --- 观察
        算法在其目标值附近如何“呼吸”。
  \item 证明交错调和级数的某种重排
        发散到 $+\infty$ \emph{（每块中取足够多的正项，
        使每次增加 $1$，利用
        第 12 问，各块之间以单个负项隔开）}。
  \item 用 $\gamma$ 律细化\cref{ex:b1:series:harmonicstack}，
        证明使 $H_N \geq
        20$ 的首个指标满足 $N = \eu^{\,20 - \gamma}\,(1 + o(1))
        \approx 2.7\cdot10^{8}$ --- 欧拉常数恰好是
        粗略估计所缺失的那一项修正。
  \item 综合各点，每点用一句话说明：（i）$\gamma$ 律及其
        三部分（$\ln n$、$\gamma$、
        $\frac{1}{2n}$）各自的贡献；（ii）为何条件收敛
        使级数的和依赖于求和次序，而绝对收敛
        则禁止这种情形；（iii）$(p,q)$ 公式如何
        是借助 $\gamma$ 律的\emph{计算}，而非
        抽象的存在性断言；（iv）这些线索
        在哪里延续 --- 本科二年级分册中的幂级数与
        级数的乘积，以及本科三年级分册中的斯特林公式
        周末问题，在那里同样的“和与积分”的
        比较手法发挥其全部威力。
\end{enumerate}
\end{problem}
